问题二以 2022 年公报口径中的江豚 1249 头和长江鲟 438 尾作为初值与口径锚定，其中 \(\delta_B\) 表示通道阻隔和栖息地碎片化的综合损失，\(R_S\) 表示人工放流强度。基线修复情景下，江豚和长江鲟的归一化末值分别为 1.133 和 1.140。仅提高通道阻隔且不叠加污染时，两者的末值为 0.901 和 0.856，较基线分别下降 20.5\% 和 24.9\%；将人工放流项置为 \(R_S=0\) 后，末值分别为 1.145 和 0.519，江豚较基线变化 +1.0\%，长江鲟下降 54.5\%。污染情景下，两者的末值分别为 0.623 和 0.706，较基线分别下降 45.0\% 和 38.0\%。

从机制上看，这些对照只支持分项比较：不叠加污染的通道阻隔情景给出阻隔作用相对基线的变化，\(R_S=0\) 的对照用于拆分自然恢复与人工放流补偿效应，污染情景则用于比较污染叠加阻隔对基线收益的侵蚀。上述情景输出都是当前参数组与 2022 年公报归一化口径下的内部投影，不构成独立外部验证；\(M\) 和 \(V\) 是传导和竞争状态，不是本问的直接观测目标。
